\selectlanguage{ngerman}%
\frenchspacing%
\begin{ledgroupsized}[r]{120mm}%
\footnotesize%
\pstart%
\noindent%
\textbf{Überlieferung:}%
\pend%
\end{ledgroupsized}%
\begin{ledgroupsized}[r]{114mm}%
\footnotesize%
\pstart%
\parindent -6mm%
\makebox[6mm][l]{\textit{L}}%
Notiz:
LH IV\;1,\;2a Bl.\;15.
Ein Zettel, %
Ränder abgerissen.
Eine Seite auf Bl.\;\href{\Scan{LH_04_01_02a_015r}}{15\:r\textsuperscript{o}}.
Bl.\;\href{\Scan{LH_04_01_02a_015v}}{15\:v\textsuperscript{o}}
ist unbeschrieben.
\pend%
\pstart%
\parindent -6mm%
\makebox[6mm][l]{\textit{E}}%
Zitat:
\textsc{F.\;Duchesneau}, \textit{Le modèles du vivant de Descartes a Leibniz}, Paris 1998, S.\;350, Anm.\;1.\cite{02588}
\pend%
\end{ledgroupsized}%
\vspace{5mm}%
\begin{ledgroup}%
\footnotesize%
\pstart%
\noindent%
\textbf{Datierungsgründe:}
Der Träger der undatierten Notiz N.\;\ref{39746}, die eine Analogie zwischen den organischen und den mechanischen Naturgesetzen hinsichtlich ihrer gemeinsamen Postulierung von \textit{sytemata in systematibus} feststellt, weist weder Spuren eines Wasserzeichens noch sonstige datierungsrelevante äußere Merkmale auf.
Die im Text verwendete Terminologie -- insbesondere der zweimal vorkommende Begriff "plastisch" -- ermöglicht jedoch eine mutmaßliche Datierung.
Es ist nämlich zu vermuten, dass die Notiz N.\;\ref{39746} im Zusammenhang mit der Debatte zwischen Pierre Bayle\indexN{\textls{Bayle} (Baylius, Beill, Baile), Pierre 1647\textendash1706} und Jean Le Clerc\indexN{\textls{Le Clerc} (Clericus, Clerique, Clerke, Liberius de Sancto Amore), Jean 1657\textendash1736} über die Existenz und Bedeutung der plastischen Naturen entstand.
Damit befasste sich Leibniz ausführlich in einem im Maiheft 1705 der \textit{Histoire des ouvrages des savans} veröffentlichten Aufsatz (\textit{Considerations sur les principes de vie, et sur les natures plastiques}, \textit{LPG} VI, S.\;539--546)\cite{02589}\cite{02039} sowie in der inhaltlich und zeitlich verbundenen Aufzeichnung \textit{Éclaircissement sur les natures plastiques et les principes de vie et de mouvement} (\textit{LPG} VI, S.\;546--555).\cite{02590}
(Siehe über diesen thematischen Zusammenhang \textsc{F.\;Duchesneau}, \textit{Le modèles du vivant de Descartes a Leibniz}, Paris 1998, S.\;349\,f.)\cite{02588}
Es ist somit plausibel anzunehmen, dass Leibniz im Zuge seiner Auseinandersetzung mit dieser Thematik auch die Notiz N.\;\ref{39746} verfasst hat, d.\,h. im gleichen Jahr 1705 oder später.
Eine frühere Datierung kann aber nach heutigem Wissensstand nicht ausgeschlossen werden.
\pend%
\end{ledgroup}%
\selectlanguage{latin}%
\frenchspacing%
\vspace{8mm}%
\normalsize%
\pstart%
\noindent%
\lbrack%
\href{\Scan{LH_04_01_02a_015r}}%
{15\:r\textsuperscript{o}}%
\rbrack\
Duplices naturae leges%
\indexS{lex naturae}%
\indexS{natura}%
\lbrack:\rbrack\
dynamicae%
\indexS{lex naturae dynamica}
et plasticae%
\indexS{lex naturae plastica}
seu organicae.%
\indexS{lex naturae organica}
Est tamen et in dynamicis%
\indexS{dynamicum}
hoc velut organicum,%
\indexS{organicum}
quod obtineri non possint
nisi materia ubique Elastica esset,%
\indexS{materia elastica ubique}
neque elastrum%
\indexS{elastrum ubique in materia}
\edtext{%
ubique%
}{%
\lemma{ubique}%
\Bfootnote{%
\textit{erg.\;L}%
}}
in materia,%
\indexS{materia elastica ubique}
nisi sytemata in systematibus collocarentur.%
\indexS{systemata in systematibus}
In quod dynamica%
\indexS{dynamicum}
respondent plasticis,%
\indexS{plasticum}
quae semper organa in organis habent.%
\indexS{organa in organis}
\pend%
\selectlanguage{latin}%
\frenchspacing%
\normalsize%
